\documentclass[preprint,12pt]{elsarticle}

\usepackage{amsmath,amssymb}
\usepackage{booktabs}
\usepackage{tabularx}
\usepackage{graphicx}
\usepackage{microtype}
\usepackage{xcolor}
\usepackage{placeins}
\usepackage{hyperref}
\hypersetup{hidelinks}

\journal{Journal of Information Security and Applications}

\newcommand{\Lang}{\mathcal{L}}
\newcommand{\Rank}{\mathsf{Rank}}
\newcommand{\Unrank}{\mathsf{Unrank}}
\newcommand{\Derive}{\mathsf{Derive}}
\newcommand{\Perm}{\Pi}
\newcolumntype{R}[1]{>{\raggedright\arraybackslash}p{#1}}

\begin{document}

\begin{frontmatter}

\title{Exact Policy-Space Credential Derivation with Factor-Preserving Root-Key Recovery}

\author{Yuanyi Zhang}
\address{Independent Researcher}

\begin{abstract}
Legacy enterprise services often retain password-only authentication, leaving
a client to store or regenerate a separate credential for each account. Exact
Policy-Space Credential Derivation (EPSCD) compiles a bounded policy into a
finite representation, counts and indexes its accepted language, and maps
credential generations through a context-separated keyed permutation before
unranking. This specializes established finite-language and rank-then-encipher
techniques to rotation. Successful outputs form a without-replacement sequence
inside one credential lineage. If $q$ prior generation/password pairs are
observed, an ideal random permutation leaves an unseen generation uniform over
the $N_P-q$ unused legal values; a real backend supplies only the corresponding
computational PRP claim. Recoverability creates a separate systems problem:
logical 2-of-3 sharing can collapse when an endpoint holding one share can
automatically invoke honest infrastructure to release another. We define a
deployment capability closure and instantiate a factor-preserving profile with
heterogeneous device, offline, and network shares plus independently signed,
session-bound network release. Freshness checkpoints bind credential, policy,
Root-Key, and share-set generations, separating ordinary re-sharing from true
Root-Key replacement. On the published SOUPS 2022 corpus, 120 of 121 exact
translations compile under one fixed budget, and warm complete derivation has
0.199~ms median latency on the test host. Adversarial tests and bounded model
checking exercise exposure, rollback, and recovery guards. The artifact remains
a research prototype with a bounded FF1 backend, no production recovery
transport, and no target-service adapters.
\end{abstract}

\begin{keyword}
deterministic password generation \sep password policy \sep finite automata
\sep format-preserving encryption \sep credential rotation \sep key recovery
\sep rollback protection
\end{keyword}

\end{frontmatter}

\section{Introduction}

Passwordless authentication is preferable when a service can be changed, but
many internal Web applications, appliances, and vendor systems expose only a
password interface. Passwords are therefore likely to persist even in managed
environments~\cite{herley2012persistence}. A client-side manager must retain or
regenerate service credentials. Published deterministic generators including
PwdHash, Password Multiplier, and PALPAS establish site-separated derivation
and non-secret per-site state as useful design points
\cite{ross2005pwdhash,halderman2005multiplier,palpas}. These published
specifications do not define the exact generation-indexed,
without-replacement sequence studied here.

Two questions are easily blurred. First, does every random choice made by an
encoder avoid modulo bias? Second, is the \emph{final} output uniform over all
strings accepted by the complete policy? The first property does not imply the
second. For example, a generator that first inserts one required character per
class and then shuffles can assign different masses to legal strings with
different class multiplicities. A second issue appears during rotation: a
deterministic function evaluated independently at generations $g_1\ne g_2$
can return the same password. The probability may be small in a large domain,
but there is no structural non-repetition guarantee.

Established work already supplies the mathematical tools needed to address
these issues. Accepted words of a finite regular language can be counted and
ranked, and rank-then-encipher maps arbitrary finite formats through a keyed
permutation~\cite{goldberg1985ranking,black2002finite,bellare2009fpe}. Our
research question is therefore not whether these tools can be reinvented for
passwords. It is whether legacy password-only accounts can be managed without
stored service-password values while preserving a mathematically defined
rotation sequence and a recovery threshold that does not collapse under
deployment capabilities.

The contributions are:

\begin{enumerate}
\item An exact, exposure-aware policy-space credential sequence. It binds
  credential generation to one finite-domain permutation, separates
  non-repetition from unpredictability, characterizes the ideal residual space
  after $q$ observed credentials, and defines credential-key rekeying and
  authenticated history semantics.
\item A factor-preserving heterogeneous Root-Key recovery profile. It defines
  deployment capability closure, exhibits logical-threshold collapse, and
  prevents an endpoint request capability from becoming authority to release a
  second top-level share.
\item An integrated rotation, recovery, compromise, and freshness lifecycle,
  supported by a Rust artifact, published-corpus and published-baseline
  measurements, adversarial tests, and bounded model checking. This is
  evaluation and systems integration, not a third cryptographic primitive.
\end{enumerate}

Counting, ranking, rank-then-encipher, FF1, HKDF, and Shamir sharing remain
standard components. Ed25519 and X25519 are also standard mechanisms. The
novelty boundary is the credential exposure and deployment-threshold
specialization, not a new primitive.

\section{Related work and novelty boundary}

\subsection{Deterministic credential generation}

PwdHash derives site-specific passwords in the browser
\cite{ross2005pwdhash}. Password Multiplier similarly combines a master secret
with a site name and supports browser-based use
\cite{halderman2005multiplier}. PALPAS derives credentials from a high-entropy
secret and per-site salts while synchronizing non-secret parameters
\cite{palpas}. These works establish that deterministic service separation and
non-secret site metadata are prior art. Their published specifications do not
state an exact accepted-language index or a generation-indexed permutation
that samples without replacement. Table~\ref{tab:published-properties} uses
``not specified'' when a cited work does not establish a property.

Grilo et al. formally verify a random password generator for composition
policies and integrate it with Bitwarden~\cite{grilo2022verified}. It provides
an important published compliance and distribution baseline, but its random
draws are not a deterministic credential sequence. EPSCD's narrower question
is how to bind an exact accepted-space index to credential generations and
rotation state without claiming that compliant or uniform password generation
itself is new.

\subsection{Policy representation and measurement}

The Password Requirements Markup Language makes service requirements
machine-readable and reports a large collection of descriptions
\cite{horsch2016prml}. Gautam et al. define a password-composition-policy
language and publish the historical 270-policy workbook used in their SOUPS
2022 study~\cite{gautam2022pcp}. Alroomi and Li later measured password-creation
policies at more than 20,000 websites~\cite{alroomi2023policies}. We use only
Gautam et al.'s public workbook as an executable corpus. The current compiler
supports a documented bounded subset and records semantic and resource
rejections; it does not infer individual policies from aggregate measurements.

\subsection{Uniform generation and finite-language indexing}

Ranking finite-language strings and uniform generation from regular languages
have a long history. Published algorithms span context-free counting,
space-efficient algebraic-language generation, controlled floating-point
approximations, linear expected-time regular-language sampling, and Dichopile's
time/space trade-off
\cite{hickey1983uniform,goldwurm1995random,denise1999uniform,
bernardi2012linear,oudinet2013uniform}. Goldberg and Sipser provide related
ranking foundations~\cite{goldberg1985ranking}. These results make exact
counting and uniform formal-language generation prior art. Dichopile is used in
the evaluation for that reason; it is not a deterministic credential-system
competitor.

\subsection{Finite-domain encryption}

Black and Rogaway formalized ciphers on arbitrary finite
domains~\cite{black2002finite}. Bellare et al. described rank-then-encipher FPE
and applied ranking and unranking directly to regular
languages~\cite{bellare2009fpe}. The composition of a regular language,
rank/unrank, and a keyed permutation is therefore established. EPSCD gives
credential generations, policy epochs, stable salt, and authenticated history
specific semantics around that composition.

NIST SP~800-38G standardizes FF1 and FF3-family format-preserving encryption;
the 2025 revision draft removes FF3 and requires an FF1 domain of at least one
million~\cite{nist38g}. Small domains also require care against message-recovery
and codebook-style attacks~\cite{durak2017fpe,hoang2018curse}. We use FF1 only as a replaceable
prototype backend and enforce the draft's minimum domain. The security
statements are parameterized by a finite-domain permutation interface rather
than equating the interface with this bounded implementation.

\subsection{Observed credentials and key exposure}

PRP security already permits known or chosen input/output pairs; conditioning
an ideal random permutation on such pairs leaves a uniform permutation over
the unused inputs and outputs. We do not claim this symmetry as new. The
contribution is to state what an exposed password means in the indexed
credential protocol and to separate the ideal residual-space statement from a
real backend's computational bound. Identity-based FPE has also studied how to
localize damage from FPE key exposure~\cite{bellare2017ibfpe}. Consequently,
per-credential key separation and rekeying are lifecycle requirements, not a
new key-exposure primitive.

\subsection{Rotation lifecycle and Root-Key recovery}

Password changes cross a boundary that a client cannot transact atomically:
the remote service may accept an update while the response is lost. EPSCD
therefore separates candidate selection, submission, observation, and local
commit, and treats positive update responses as evidence to be verified rather
than as a commit certificate. This is a protocol specialization around the
credential sequence, not a claim to a new distributed-transaction primitive.

Shamir secret sharing provides the standard threshold mechanism used by the
2-of-3 Root-Key recovery profile~\cite{shamir1979share}. MFKDF is
published prior work on multi-factor derivation and threshold client-side
recovery~\cite{nair2023mfkdf}; subsequent cryptanalysis shows why state
integrity and precisely stated assumptions matter in recovery constructions
\cite{scarlata2024mfkdf}. Threshold OPRF-based password-protected sharing is
also established~\cite{jarecki2017toppss}.

Distributed recovery systems constrain the systems claim. SafetyPin distributes
PIN-protected backup recovery across HSMs and logs and punctures recovery
capability~\cite{dauterman2020safetypin}. SVR3 splits trust across
heterogeneous enclave types and cloud providers and supplies rollback
protection~\cite{connell2024svr3}. Little et al. give a deployed, privacy-
preserving server-assisted account-recovery protocol~\cite{little2024recovery},
and SPHINX studies password management under device compromise
\cite{shirvanian2019sphinx}. Thus heterogeneous, threshold, authorized, and
fresh recovery are individually prior art. Our narrower question is whether an
endpoint already holding one top-level share can invoke honest infrastructure
to obtain another. The proposed closure analysis and lifecycle profile target
that deployment-threshold failure; they do not claim a new sharing or public-
key primitive.

\begin{table*}[t]
\centering
\scriptsize
\setlength{\tabcolsep}{2.0pt}
\caption{Comparison of published work by stated generation semantics. ``Not
specified'' means that the cited paper does not establish the property; it is
not a claim that an extension is impossible.}
\label{tab:published-properties}
\resizebox{\textwidth}{!}{%
\begin{tabular}{@{}R{0.105\textwidth}R{0.085\textwidth}R{0.14\textwidth}R{0.13\textwidth}R{0.14\textwidth}R{0.16\textwidth}R{0.14\textwidth}@{}}
\toprule
Method & Status & Generation model & Policy model & Space characterization & Sequence semantics & Rotation semantics \\
\midrule
PwdHash \cite{ross2005pwdhash} & USENIX Security 2005 & site-scoped hash derivation & output transformation & not specified & one site-scoped value & not specified \\
Password Multiplier~\cite{halderman2005multiplier} & WWW 2005 & site-scoped hash derivation & not specified & not specified & one site-scoped value & not specified \\
PALPAS \cite{palpas} & ARES 2015 & secret and per-site salt & server parameters & not specified & one value per salt state & salt-state update \\
Verified RPG~\cite{grilo2022verified} & iFM 2022 & random sampling & composition policy & accepted distribution specified & independent random draws & not applicable \\
Dichopile \cite{oudinet2013uniform} & TCS 2013 & uniform random sampling & regular-language automaton & path-count weights & independent random draws & not applicable \\
EPSCD & this work & generation-indexed keyed permutation & bounded policy IR & exact indexed language & without replacement within an epoch & authenticated history and remote evidence \\
\bottomrule
\end{tabular}
}
\end{table*}


\section{Model and security assumptions}

The client serves legacy applications that publish a finite password policy
and accept a username and password. Public information includes $P$, its
accepted language $\Lang_P$, $N_P$, rank/unrank, scheme and encoder versions,
vault/service/account/policy identifiers, the credential salt, policy hash,
policy and Root-Key generations, credential generation, share-set generation,
and authenticated history descriptors. Secrecy never depends on hiding this
metadata. The derivation secrets are $K_{root}$ and per-credential $K_{cred}$;
service-password plaintext exists only during derivation and use.

We consider exposure of one or $q$ generation/password pairs, a compromised
server verifier, compromise of $K_{cred}$ or $K_{root}$, replay of a complete
old authenticated snapshot, and compromise of recovery domains. Server
verifier compromise permits direct enumeration of $\Lang_P$ against the
verifier and bypasses the Root Key and permutation. Thus a 256-bit Root Key
does not imply 256-bit service-password guessing strength.

Recovery uses four administrative or physical domains: D is a managed
endpoint, U a cloneable removable/offline medium, N the network recovery
infrastructure, and A an independent release-authorization domain. A holds no
Root-Key share. For a compromised domain $X$, $\mathsf{Closure}(X)$ contains
the files, accessible secure-storage contents, TLS and API credentials,
cookies, request-signing keys, automatic protocol calls, and responses that
honest services release to those capabilities. The intended top-level 2-of-3
access structure is factor-preserving only if
\[
|\mathsf{Closure}(X)\cap\{S_D,S_U,S_N\}|<2
\]
for every single protected compromise domain $X$. Independence is a deployment
assumption: one administrative account controlling D, N, and A invalidates it.

We assume HKDF behaves as a PRF, the finite-domain permutation meets its stated
query and domain bounds, Ed25519 signatures are unforgeable, X25519-based
session encryption protects released fragments, and authenticated storage plus
an independently administered checkpoint supplies freshness. A fully
compromised endpoint is assumed to obtain secrets accessible while it is
unlocked; platform secure storage is hardening, not an exception. Online
guessing, malicious policy publication, administrator social engineering,
denial of service, and target-service availability remain outside the claimed
guarantees.

\begin{figure}[t]
\centering
\small
\begin{tabular}{R{0.27\linewidth}R{0.63\linewidth}}
\toprule
Class & Contents \\
\midrule
Public & $P,\Lang_P,N_P,\Rank,\Unrank$, identities, salts, generations, hashes, scheme versions \\
Derivation secrets & $K_{root},K_{cred}$ and plaintext credentials during use \\
Recovery secrets & top-level shares $S_D,S_U,S_N$ and N's 3-of-5 fragments \\
Authorization & A-domain Ed25519 private keys; no Root-Key share \\
Freshness & public monotonic checkpoint over Root, share set, policy, lineage, and credential generation \\
\bottomrule
\end{tabular}
\caption{Exposure classes. Secrecy does not depend on hiding the policy
language or derivation metadata.}
\label{fig:exposure-boundary}
\end{figure}


\section{Policy model and exact bijection}

\subsection{Supported policy IR}

A policy $P$ defines a finite language $\Lang_P\subseteq\Sigma^*$ with bounded
length $L_{min}\le |w|\le L_{max}$. The version-1 policy IR supports an ordered
effective alphabet, forbidden characters, class minimum and maximum counts,
fixed positions, fixed prefix and suffix, forbidden first and last characters,
per-character maximum totals, maximum identical and sequential ASCII runs, and
finite forbidden substrings. `null', empty, and absent fields are distinguished
by canonical JSON serialization. The alphabet order is part of the protocol.

The compiler constructs a reachable deterministic product state containing
only data required by enabled constraints. Minimum counters saturate after the
minimum is met; maximum counters do not. Prefix/suffix and fixed-position
constraints prune transitions. A bounded suffix state recognizes forbidden
substrings. The compiler rejects an empty language, inconsistent class,
unsupported normalization, length above 128, or more than 250,000 reachable
states. Thus ``supported'' means expressible by this IR \emph{and} compilable
within its explicit resource budget, not arbitrary PCRE or semantic policy
code.

\subsection{Exact counting}

Let the compiled automaton be
$A_P=(Q,\Sigma,\delta,q_0,F,L_{min},L_{max})$. For a target length $\ell$, define

\begin{equation}
C_\ell(i,q)=\left|\{x\in\Sigma^{\ell-i}:\delta^*(q,x)\in F\}\right|.
\end{equation}

The boundary and recurrence are

\begin{align}
C_\ell(\ell,q)&=\begin{cases}1&q\in F\\0&q\notin F,\end{cases}\\
C_\ell(i,q)&=\sum_{c\in\Sigma:\delta(q,c)\ne\bot}
 C_\ell(i+1,\delta(q,c)).
\end{align}

The exact space is
$N_P=\sum_{\ell=L_{min}}^{L_{max}}C_\ell(0,q_0)$. Every count is an arbitrary-
precision non-negative integer. `entropyBits' reports $\log_2 N_P$ as the
entropy of a uniform random variable over the language; a fixed keyed
derivation output is deterministic and must not itself be described as having
that entropy.

\subsection{Rank and unrank}

Strings use shortlex order: length first, then the canonical alphabet order.
To unrank $r\in[0,N_P)$, the algorithm subtracts complete length blocks, then
examines each outgoing character in order. The count on that edge is the size
of its contiguous rank interval. Ranking performs the inverse additions.

\textbf{Lemma 1 (bijection).} If the count table is correct,
$\Unrank_P:[0,N_P)\leftrightarrow\Lang_P$ is a bijection and $\Rank_P$ is its
inverse.

\textit{Proof sketch.} At each state, the recurrence partitions the completion
set into disjoint, adjacent intervals indexed by the next character. Their
sizes sum to the current completion count. Induction on the remaining length
gives existence and uniqueness; the outer shortlex length blocks are disjoint
by construction.\hfill$\square$

\section{Generation-indexed finite-domain derivation}

Figure~\ref{fig:pipeline} gives the scheme-version-1 data flow. A 256-bit
credential key is derived from the Root Key with HKDF-SHA256, using a 128-bit
credential salt and an RFC-8785 canonical context:

\begin{figure*}[t]
\centering
\small
\setlength{\tabcolsep}{5pt}
\resizebox{\textwidth}{!}{%
\begin{tabular}{ccccccccc}
\fbox{Policy $P$} & $\rightarrow$ & \fbox{Canonical IR} & $\rightarrow$ & \fbox{Finite automaton} & $\rightarrow$ & \fbox{Exact counts} & $\rightarrow$ & $N_P,\Rank_P,\Unrank_P$ \\
&&&&&&&& $\downarrow$ \\
\fbox{$K_{root}$ + context $C$} & $\rightarrow$ & \fbox{$K_{cred}$ + tweak $T$} & $\rightarrow$ & \fbox{generation $g$} & $\rightarrow$ & \fbox{$r=\Perm_{K,T}(g)$} & $\rightarrow$ & \fbox{$\Unrank_P(r)$} \\
\end{tabular}
}
\caption{EPSCD data flow. Policy compilation defines one exact finite domain; credential generation is the input to one context-bound permutation, not part of its tweak.}
\label{fig:pipeline}
\end{figure*}


\begin{align}
C_K &= \mathsf{JCS}(\texttt{EPSCD/credential-key/scheme-v1},\nonumber\\
&\quad schemeVersion,vaultID,serviceID,accountID,\nonumber\\
&\quad rootGeneration),\\
K_{cred} &= \mathsf{HKDF\mbox{-}SHA256}(K_{root},credentialSalt,C_K).
\end{align}

The exact byte-level construction is specified by canonical JSON contexts and
the artifact's fixed vector. The public `schemeVersion' is 1. The permutation
tweak binds its domain label, scheme version, vault, service, account,
credential salt, Root-Key generation, policy identity, policy version, policy
hash, and policy epoch.
Credential generation is deliberately excluded from the tweak because it is
the permutation input.

Let $\Perm_{K,T}$ be a total keyed permutation on $[0,N_P)$. For generation
$0\le g<N_P$,

\begin{equation}
\Derive(K,T,g,P)=\Unrank_P(\Perm_{K,T}(g)).
\end{equation}

The prototype implements the interface using FF1-AES-256 over the smallest
binary superset followed by cycle walking. It rejects domains below one
million, domains above 512 bits, inputs outside the domain, and derivations that
exceed 1,024 walks. This limit makes the backend a partial, failure-closed
implementation. The following properties therefore refer to the abstract
total permutation, or to successful calls of the prototype where stated.

Let $M=2^{\lceil\log_2N_P\rceil}$. Since $N_P>M/2$ for a non-power-of-two
domain, a permutation orbit can visit at most $M-N_P$ outside points before
returning to the accepted subdomain. Thus unbounded cycle walking terminates
within $M-N_P+1$ primitive calls. For an ideal random permutation,
\begin{equation}
\Pr[W>k]=\frac{(M-N_P)_k}{(M)_k}
\le\left(\frac{M-N_P}{M}\right)^k<2^{-k},
\end{equation}
where $(x)_k$ is a falling factorial and the numerator is zero for
$k>M-N_P$. If $N_P=M$, one primitive call always returns to the domain. A
1,024-call cap fails in the ideal model with probability $\Pr[W>1024]$, below
$2^{-1024}$ in the non-power-of-two case. This calculation is not an
unconditional availability guarantee for concrete FF1; the full derivation is
given in Appendix~\ref{app:formal-proofs}.

\textbf{Theorem 1 (compliance).} Every successful derivation belongs to
$\Lang_P$. This follows directly from Lemma~1.

\textbf{Theorem 2 (epoch-local non-repetition).} For fixed $(K,T,P)$ and
$g_1\ne g_2$ in the domain, successful outputs differ. Both $\Perm$ and
$\Unrank$ are injective.

\textbf{Theorem 3 (ideal-permutation marginal).} If $\Perm$ is sampled
uniformly from all permutations on $[0,N_P)$, then for fixed $g$ and any
$w\in\Lang_P$,
$\Pr[\Derive(K,T,g,P)=w]=1/N_P$. A real PRP gives computational
indistinguishability under its security assumptions, not unconditional exact
uniformity over keys.

If an adversary observes $password_g$, public ranking reveals the pair
$(g,\Rank_P(password_g))$. Injectivity prevents this rank from reappearing at a
different input, but says nothing about predicting the remaining mapping:
injectivity provides non-repetition; PRP security provides unpredictability
under observed input/output pairs.

\textbf{Theorem 4 (observed-credential residual uniformity).} Condition an
ideal random permutation on $q<N_P$ distinct, consistent generation/rank pairs.
For any unseen generation, each unused rank has conditional probability
$1/(N_P-q)$. Equivalently, the derived password is uniform over the
$N_P-q$ unobserved legal values.

There are $(N_P-q)!$ permutations consistent with the observations and
$(N_P-q-1)!$ after fixing one additional unseen pair. Appendix
\ref{app:formal-proofs} gives the complete proof and a game hop: a real
backend replaces this equality with computational indistinguishability bounded
by its PRP advantage. This standard random-permutation symmetry is specialized
to the credential sequence; it is not a new permutation theorem and does not
prove concrete FF1 key-recovery resistance.

\textbf{Theorem 5 (context separation).} Under HKDF and tweakable-PRP
assumptions, distinct credential identities, salts, Root-Key generations, or
policy epochs induce computationally separated derivation contexts. This is a
computational claim; outputs from different generations in one permutation are
sampling without replacement and are not statistically independent.

The key formula also fixes compromise semantics. `policyEpoch' is only a public
tweak component, so changing it cannot repair a leaked $K_{cred}$. The minimum
independent rekey starts a new credential lineage with a fresh 128-bit
credential salt and applies authenticated history exclusion before committing
its first password. Re-randomizing Shamir shares likewise cannot repair a
compromised $K_{root}$; that event requires a random replacement Root Key and
remote rotation of affected credentials.

\section{Credential sequence and rotation lifecycle}

\subsection{Protocol and state versioning}

The public protocol begins at \texttt{schemeVersion=1}. Its authenticated
credential descriptor keeps five orthogonal counters or identifiers:
\texttt{policyVersion} identifies the policy serialization and compiler
semantics; \texttt{policyEpoch} identifies one accepted-language lineage;
\texttt{rootGeneration} changes when the Root Key is replaced;
\texttt{credentialGeneration} is the permutation input within an epoch; and
\texttt{shareSetGeneration} identifies the current recovery-share set. A
same-policy rotation increments only \texttt{credentialGeneration}. A policy
change increments \texttt{policyEpoch} and resets
\texttt{credentialGeneration} to zero. Root-Key replacement increments
\texttt{rootGeneration} and installs a new share-set generation. These fields
exclusively describe the public scheme defined in this paper.

The credential salt is stable within a credential lineage. Generating a fresh
salt at every rotation would select unrelated permutations and destroy
Theorem~2's sequence interpretation. A new salt explicitly creates a new
lineage, so its first value must still pass authenticated history exclusion.

\subsection{History exclusion across policy epochs}

Independent policy epochs can have intersecting languages, so Theorem~2 does
not prevent cross-epoch reuse. For a server history window $h$, the client
retains authenticated descriptors of the recent committed versions and
rederives their passwords; it does not store the values. Duplicate values are
removed before forming the exclusion set $E$. The pending epoch tests
successive unused generation inputs and selects the first output outside $E$.

\textbf{Proposition 1 (bounded exclusion).} Under a total permutation, if at
least $|E\cap\Lang_P|+1$ unused generation inputs remain, the process terminates
after at most $|E\cap\Lang_P|+1$ candidates. It fails when history cannot be
authenticated or rederived, including history under an unavailable Root-Key
generation. A configured history window above 64 is rejected.

\textbf{Theorem 6 (uniform exclusion marginal).} Let $e=|E\cap\Lang_P|<N_P$.
Under an ideal random permutation, scanning distinct generation inputs and
returning the first output outside $E$ produces every
$w\in\Lang_P\setminus E$ with probability $1/(N_P-e)$. It uses at most $e+1$
attempts and has expected attempt count $(N_P+1)/(N_P-e+1)$.

The proof maps the deduplicated exclusion set to ranks and observes that the
images of distinct scanned generations form an ordered sample without
replacement. Symmetry gives the marginal, while the minimum order statistic of
the $N_P-e$ admissible ranks gives the expectation. Appendix
\ref{app:formal-proofs} supplies the complete argument. The theorem fixes an
authenticated local history independently of the ideal permutation; it says
nothing about undocumented server state.

The client-side check cannot predict undocumented server history rules. Such a
server rejection is outside the policy-space algorithm and must not silently
change the committed generation metadata.

\subsection{Credential-space adequacy}

Let $B_P=\log_2N_P$. This is the exact entropy of an ideal uniform draw from
the accepted language, not the entropy of a fixed deterministic output. Two
deployment predicates must remain separate. `backendDomainEligible' asks
whether the concrete finite-domain implementation accepts $N_P$; the prototype
requires at least one million values and at most 512 bits. A configurable
`credentialSpaceFloor' instead asks whether $B_P$ is adequate for the service's
online/offline guessing threat profile. Satisfying the first does not imply the
second. If an attacker obtains the server verifier, salt, and hash parameters,
it can enumerate $\Lang_P$ directly, regardless of Root-Key length.

\subsection{Evidence-bounded remote rotation}

Figure~\ref{fig:rotation-lifecycle} separates derivation from remote evidence.
From \textsc{Active}, the client selects a history-safe generation and durably
records it as \textsc{Pending} before calling the target adapter. Submission
enters \textsc{Submitted}. A lost response enters
\textsc{UnknownOutcome}; neither state permits local commit. The adapter then
probes authentication with both the candidate and currently committed
password, subject to a lockout budget and retry backoff.

\begin{figure*}[t]
\centering
\resizebox{\textwidth}{!}{$
\fbox{\textsc{Active}$(g)$}
\xrightarrow{\text{persist }(g,g+1,opID)}
\fbox{\textsc{Prepared}}
\xrightarrow{\text{submit}}
\fbox{\textsc{Submitted}}
\xrightarrow{\text{adapter observations}}
\fbox{\textsc{Verifying}}
$}

\vspace{6pt}
\scriptsize
\begin{tabularx}{0.96\textwidth}{>{\centering\arraybackslash}X>{\centering\arraybackslash}X>{\centering\arraybackslash}X}
$\mathsf{CommitEvidence}=\mathsf{true}$
$\rightarrow\textsc{Committed}(g+1)$ &
Conclusive old-authoritative evidence
$\rightarrow\textsc{Aborted}(g)$ &
Insufficient or contradictory evidence
$\rightarrow\textsc{UnknownOutcome}$ \\
\multicolumn{3}{c}{HTTP 200 $\not\Rightarrow$ commit; \quad
\textsc{UnknownOutcome} preserves reconstructibility of both $g$ and $g+1$.}
\end{tabularx}
\caption{Evidence-bounded remote rotation. Persisting a candidate and sending
an update never commits local metadata. The adapter contract determines whether
new acceptance, new-only acceptance, atomic evidence, or authoritative version
readback is sufficient.}
\label{fig:rotation-lifecycle}
\end{figure*}


The four observations have distinct outcomes. \textsc{NewOnly} permits an
authenticated local commit; \textsc{OldOnly} permits abort of the candidate.
\textsc{Both} is an ambiguous remote state, for example during replicated
authentication transition, and \textsc{Neither} may indicate a third value,
lockout, or adapter failure. Both enter \textsc{Recovering} and require explicit
reconciliation. A successful update response without these probes is not
sufficient evidence.

Each target adapter implements
\texttt{submit\_update}, \texttt{classify\_response},
\texttt{verify\_new}, \texttt{verify\_old}, a lockout budget, and retry
backoff. Verification semantics are target-specific: an HTTP status or success
page alone need not demonstrate that subsequent authentication observes the
new value. The prototype specifies this contract and models its state
transitions, but does not yet include production adapters.

\section{Factor-preserving Root-Key recovery}

The Root Key is a uniformly sampled 256-bit value. Standard Shamir sharing
creates $(S_D,S_U,S_N)$ with threshold two. D retains $S_D$, U retains $S_U$,
and $S_N$ is split again into 3-of-5 network fragments. This second sharing
changes N's availability and compromise threshold but does not create another
top-level Root-Key factor.

\begin{figure*}[t]
\centering
\small
\resizebox{\textwidth}{!}{%
\begin{tabular}{ccccc}
\fbox{D: managed endpoint $S_D$} & & \fbox{U: removable/offline $S_U$} & &
\fbox{N: 3-of-5 fragments of $S_N$} \\
& \multicolumn{3}{c}{$\searrow\qquad \text{Shamir 2-of-3}\qquad\swarrow$} & \\
& \multicolumn{3}{c}{\fbox{$K_{root}$}} & \\
& \multicolumn{3}{c}{$\downarrow\ \mathsf{HKDF}(\text{credential identity and lineage})$} & \\
& \multicolumn{3}{c}{\fbox{$K_{cred}$ $\rightarrow$ EPSCD $\rightarrow$ service passwords}} & \\
\multicolumn{5}{c}{A: independent approval signs a session-bound release ticket; A holds no Root-Key share.}
\end{tabular}}
\caption{Unified derivation and recovery architecture. Recovery reconstructs
only the random Root Key; service-password values are never recovery objects.}
\label{fig:unified-architecture}
\end{figure*}

\FloatBarrier

\subsection{Operational profile}

In the intended day-to-day profile, no complete Root Key is persistently
stored. D retains a platform-protected $S_D$, while cloneable removable medium
U stores an authenticated, generation-bound $S_U$ envelope. The current
artifact does not password-encrypt that USB envelope; a deployment may add
password-based authenticated encryption as hardening, in which case the
password unlocks the envelope rather than generating the share and does not
constitute an additional Shamir factor. To create, display, or rotate a
credential, the user connects U and makes $S_U$ available; $S_D$ and $S_U$ are
combined to reconstruct $K_{root}$ transiently in the client process, from
which $K_{cred}$ and the EPSCD output are derived. Mutable secret buffers are
cleared on a best-effort basis after use, which is hygiene rather than a memory-
confidentiality guarantee. $S_N$ is reserved for an alternative recovery path
when D or U is unavailable. Thus neither service-password values nor a complete
Root Key are persistent recovery objects, but a fully compromised endpoint can
observe them during a legitimate unlocked reconstruction session.

\subsection{Threshold collapse and closure}

A naive endpoint can store both $S_D$ and a credential that honest network
nodes automatically honour to release $S_N$. Then
$S_D,S_N\in\mathsf{Closure}(D)$ and compromise of one physical domain recovers
the Root Key. Splitting $S_N$ among five nodes does not fix the defect when the
same endpoint can invoke three honest nodes. Cryptographic threshold and
deployment compromise threshold are different quantities.

\begin{figure*}[t]
\centering
\small
\begin{tabular}{R{0.46\textwidth}R{0.46\textwidth}}
\toprule
Naive deployment & Factor-preserving deployment \\
\midrule
D stores $S_D$ and an automatically honoured network-release credential.
Honest nodes release $S_N$, so both $S_D,S_N\in\mathsf{Closure}(D)$ and one
endpoint compromise reconstructs $K_{root}$. &
D stores $S_D$ and request capability only. A signs the complete ticket from an
independently administered domain; three N nodes release fragments only to the
ticket-bound ephemeral public key. Thus compromise of D alone yields one
top-level share. \\
\bottomrule
\end{tabular}
\caption{Logical threshold versus deployment threshold. Threshold storage of
$S_N$ is insufficient when one endpoint can automatically invoke enough honest
nodes; independent release authorization removes that authority from
$\mathsf{Closure}(D)$.}
\label{fig:factor-collapse}
\end{figure*}


\subsection{Session-bound independent release}

D may create a request but cannot authorize release. A recovery session
contains the vault and Root-Key identifiers, \texttt{shareSetGeneration}, a
random \texttt{opID}, purpose, issue/expiry times, the authorized node set, and
an ephemeral X25519 public key $pk_R$. At least two independent A-domain
Ed25519 signers approve the RFC-8785 canonical ticket. A network node releases
its fragment only when the signatures, vault, Root-Key and share-set
generations, node authorization, lifetime, and unused/idempotent \texttt{opID}
all verify. The response is encrypted to $pk_R$ with the ticket and node
identity as associated data. Encryption protects transport but never
substitutes for authorization.

U is assumed cloneable and carries no A-domain signing secret. D carries only
request capability, not an approval key. Consequently, compromise of D alone,
copying U alone, compromising fewer than three N nodes, or compromising one A
signer supplies fewer than two top-level shares under the stated independence
assumption.

\section{Integrated recovery and freshness lifecycle}

Ordinary loss recovery reconstructs the same Root Key, increments
\texttt{shareSetGeneration}, re-shares all top-level factors, re-splits $S_N$,
and makes old tickets and fragments stale. \texttt{rootGeneration} remains
unchanged. Suspected exposure of any copied share follows the same re-sharing
path. In contrast, confirmed $K_{root}$ compromise samples a new Root Key,
increments both Root-Key and share-set generations, derives fresh credential
keys, and requires evidence-bounded remote rotation of every affected account.

Freshness is public but externally anchored. The prototype's persistent CAS
checkpoint covers Root-Key and share-set generations and, per credential,
policy epoch, credential generation, and a public credential-lineage digest.
A complete old authenticated snapshot is rejected when any component precedes
the checkpoint. This supplies the deployment condition omitted by Theorem~2:
no-repeat requires both permutation injectivity and generation freshness.
Freshness-service credentials belong to $\mathsf{Closure}(D)$ but cannot sign a
network-share release ticket.

\section{Implementation}

The Rust prototype separates policy compilation from finite-domain derivation. The
policy compiler creates only reachable product states and stores all completion
counts as `BigUint'. Public operations include exact count by length, total
count, rank, unrank, and policy hash. Property tests compare enumeration and
counting, exercise `Rank(Unrank(r))=r', test values above 128 bits, and verify
constraint families.

`DomainPermutation' isolates the mathematical construction from the FF1
research backend. The scheme-version-1 fixed vector binds credential context,
salt, policy, Root Key, credential key, exact domain size, rank, and output
password. Callers may cache a compiled policy; complete derivation otherwise
compiles on demand. The history selector rederives the authenticated exclusion
window and consumes successive generation inputs until it finds the first
non-excluded output. Suspected credential-key exposure invokes a distinct
rekey operation that replaces the credential salt; a policy-epoch increment
alone is not treated as remediation.

The recovery prototype uses the same manifest-selected top-level Shamir share
set as local recovery. It applies `vsss-rs' again to the canonical network
`ShareEnvelope' to obtain 3-of-5 fragments. Ed25519 approvals cover the full
canonical ticket. Each node derives an AES-GCM response key from ephemeral
X25519 and HKDF and records \texttt{(nodeID,opID,ticketHash)} for idempotency.
No service-password value, view key, data key, threshold OPRF, or opaque-object
index exists in this path.

The SQLite CAS freshness prototype stores one vault checkpoint with explicit
Root-Key and share-set generations and per-credential policy, generation, and
lineage state. Updates must dominate the previous checkpoint; reads below it
are rollback, while equal counters with a different lineage or digest are a
fork.

The current implementation limitation is that an arbitrary finite-domain PRP
backend with a proof-matched, total implementation has not replaced bounded
FF1 cycle walking. Complex forbidden-substring compilation is also a simple
suffix-state construction rather than a minimized Aho--Corasick product, which
explains one stress-policy rejection.

\section{Evaluation}

Measurements use a release build on one macOS x86-64 host. Raw JSON records
sample counts, exact integers, seeds, resource limits, and the operation timed.
The evaluation uses the published Dichopile uniform-language algorithm as an
executable baseline and otherwise reports mathematical controls or tests of
our own mechanisms. No unpublished system, local predecessor, or preprint is
used as an experimental comparison.

\subsection{Policy-corpus coverage and exact spaces}

The public SOUPS 2022 workbook contains 270 historical website policies. The
converter records the workbook digest, source row, original PCP JSON,
translation result, and rejection reason. It exactly represents 121 records in
the current IR. Each translated record is then compiled in an isolated worker
under the same predeclared budget: 250,000 reachable states, 4~GiB sampled
resident memory, and 60~s wall time. There is no length prefilter.

Compilation completes for 120 records. One policy (source row 31, length
8--50) reaches the wall-time limit with a sampled resident set near 2.0~GiB;
no worker reaches the state or memory limit. Among completed cases, the median
state count is 65 and P95 is 284. Median compilation time is 55.90~ms, while
P95 is 904.38~ms. The longest completed policies have maximum length 128.
Table~\ref{tab:realpolicy} gives the remaining distributions.

\begin{table}[t]
\centering
\small
\caption{SOUPS 2022 PCP corpus translation and compilation under a per-policy
budget of 250,000 states, 4~GiB resident memory, and 60~s wall time.}
\label{tab:realpolicy}
\begin{tabular}{lr}
\toprule
Outcome & Policies \\
\midrule
Historical source records & 270 \\
Unsupported semantic translation & 149 \\
Exact translations attempted & 121 \\
Compilation completed & 120 \\
Wall-time limit & 1 \\
State or memory limit & 0 \\
\bottomrule
\end{tabular}
\vspace{3pt}

\begin{tabular}{lrrr}
\toprule
Metric over 120 completions & Median & P95 & Maximum \\
\midrule
Reachable states & 65 & 284 & 23,227 \\
Compile time (ms) & 55.90 & 904.38 & 9,427.74 \\
Count payload (KiB) & 493.4 & 19,201.6 & 75,590.4 \\
Sampled peak RSS (MiB) & 5.6 & 55.7 & 381.6 \\
Rank median ($\mu$s) & 24.93 & 109.30 & 137.03 \\
Unrank median ($\mu$s) & 37.34 & 137.65 & 163.88 \\
\bottomrule
\end{tabular}
\end{table}


Reclassifying the 120 completed exact counts without recompilation finds that
96 satisfy the prototype backend's domain bounds. All 120 exceed illustrative
40- and 60-bit credential-space floors, while 116 exceed 80 bits. These floors
are example threat profiles, not universal recommendations; the artifact
accepts a configurable threshold and reports backend and security eligibility
separately.

The separate synthetic corpus contains 12 stress policies. Common class
constraints require 9--254 states; edge and run constraints require
6,201--8,889 states; per-character totals reach 159,250 states. The
forbidden-substring stress case reaches the 250,000-state limit.

\begin{table}[t]
\centering
\small
\caption{Selected synthetic policy compilation results (release build).}
\label{tab:policy}
\begin{tabular}{lrrr}
\toprule
Policy & States & Compile ($\mu$s) & $\log_2 N_P$ \\
\midrule
Lower/digit, length 8 & 9 & 1,703 & 41.36 \\
Four classes, length 12 & 164 & 8,891 & 71.36 \\
Four classes, length 18 & 254 & 15,912 & 111.60 \\
No symbol at edges, length 16 & 7,531 & 434,273 & 95.63 \\
Per-character total $\le2$ & 159,250 & 1,690,696 & 34.91 \\
Forbidden-substring stress & rejected & 880,798 & --- \\
\bottomrule
\end{tabular}
\end{table}


For the 18-character, 74-symbol case, the exact count is
\begin{equation*}
N_P=3923998958022941748259805567385600,
\end{equation*}
corresponding to 111.596 bits for a uniform variable over that language.

\subsection{Published uniform-generation baseline}

Oudinet et al.'s published Dichopile algorithm is the executable
regular-language baseline. We transcribed Algorithm~1 after reviewing the full
paper and use exact BigUint arithmetic rather than the floating-point mode of
its large-automata experiments. The EPSCD row uses a test-only Fisher--Yates
finite-domain permutation so that the complete 11-word language
$\{a,b0,\ldots,b9\}$ can be measured; the FF1 backend deliberately rejects
domains below one million. Each row has 100,000 samples and expected cell
counts above 9,000, satisfying the usual count condition for the Pearson
statistic.

\begin{table}[t]
\centering
\small
\setlength{\tabcolsep}{3pt}
\caption{Published-baseline and EPSCD observations on the same 11-word language
(100,000 samples per row). The table is not a system-performance ranking.}
\label{tab:dist}
\begin{tabular}{lR{0.24\linewidth}rrr}
\toprule
Method & Role & TVD & $\chi^2_{10}$ & max/min \\
\midrule
Dichopile exact~\cite{oudinet2013uniform} & published uniform-language baseline & 0.00399 & 10.24 & 1.036 \\
EPSCD test permutation & proposed sequence composition & 0.00234 & 3.59 & 1.022 \\
\bottomrule
\end{tabular}
\end{table}


Dichopile and the EPSCD test instance have TVD 0.00399 and 0.00234,
respectively; the difference is consistent with sampling variation. The EPSCD
row checks rank/permutation/unrank composition on an enumerable domain but does
not test the concrete FF1 PRP assumption. Uniform-rank and with-replacement
controls are confined to the supplement.

\subsection{Observed-credential exposure check}

An enumerable 11-rank ideal experiment conditions on $q=0,\ldots,5$ fixed
consistent generation/rank pairs and repeatedly samples the remaining
permutation. For every $q$, observed outputs reappear zero times and the
measured support is exactly $11-q$. Across 100,000 samples per $q$, TVD from
the uniform distribution on that support ranges from 0.00213 to 0.00666. This
checks the implementation of the conditional-support experiment; it neither
validates FF1 nor replaces Theorem~4.

\subsection{Recovery and rollback tests}

Adversarial tests reject endpoint-only and copied-U closures, one approval,
fewer than three node responses, expiry, public-key rebinding, changed-ticket
reuse of an \texttt{opID}, stale share-set generations, and mixed old/new
fragments. They also verify that every one of the ten 3-of-5 node combinations
reconstructs the same authenticated $S_N$, and that ordinary re-sharing makes
the old set stale. Freshness tests restore a valid generation-5-equivalent
checkpoint after later advancement and separately cover policy, Root-Key,
share-set, and credential-lineage rollback/fork cases.

In 100 release-mode local iterations, two Ed25519 approvals take a mean
369.42~$\mu$s and three node releases plus decryption and Shamir reconstruction
take 6.036~ms. These are local cryptographic baselines. The prototype has no
network transport; network RTT and human approval latency are explicitly
unmeasured and must not be inferred from these numbers.

\subsection{Local performance}

For the 18-character policy, Table~\ref{tab:performance} separates cold and
warm paths over 500 operations. Cold EPSCD starts from the policy IR, constructs
the DFA and full completion table, then derives one credential. Cold Dichopile
starts from the same IR, constructs the shared DFA, initializes length weights,
and generates one word. Warm paths reuse their respective structures. Warm
EPSCD includes HKDF, tweak serialization, FF1 cycle walking, and unrank; its
median is 198.52~$\mu$s. The exact-arithmetic Dichopile reproduction is an
operational reference for this artifact, not a benchmark against the optimized
floating-point configuration reported by Oudinet et al. The two algorithms
also produce different sequence semantics, so no speed-up factor is inferred.

\begin{table}[t]
\centering
\small
\caption{Cold and warm local release-mode timing on the same 18-character
policy (500 operations). Cold rows begin from the policy IR; warm rows reuse
their method-specific compiled structures.}
\label{tab:performance}
\resizebox{\linewidth}{!}{%
\begin{tabular}{lrrrr}
\toprule
Operation & Median & P95 & P99 & SD \\
\midrule
EPSCD cold ($\mu$s) & 7,806.13 & 8,507.33 & 9,566.33 & 754.93 \\
EPSCD warm ($\mu$s) & 198.52 & 332.16 & 425.16 & 55.94 \\
Dichopile cold ($\mu$s) & 13,793.79 & 15,084.76 & 17,140.62 & 2,609.42 \\
Dichopile warm ($\mu$s) & 5,626.27 & 6,253.10 & 7,005.11 & 567.89 \\
Rank, cached ($\mu$s) & 3.07 & 3.89 & 5.23 & 0.46 \\
Unrank, cached ($\mu$s) & 6.86 & 10.25 & 16.02 & 2.07 \\
\bottomrule
\end{tabular}%
}
\end{table}


\subsection{Cycle-walk observations}

Cycle walks were measured for every completed corpus policy within the
concrete backend's domain range. Ninety-six policies meet the one-million lower
bound and 512-bit upper bound; 24 completed policies exceed 512 bits and cannot
be instantiated by this backend. For the measured policies, 32 generations
per policy give 3,072 observations. The aggregate mean is 1.505 primitive
calls, the median is one, P95 is three, P99 is five, and the maximum is nine.
No observation reaches the 1,024-call cap.

Domain density $N/M$ ranges from 0.516 to 0.973. Figure
\ref{fig:walk-density} plots each policy mean against density and includes the
large-domain ideal reference $1/(N/M)$; the exact ideal expectation is
$(M+1)/(N+1)$. The trend supports the dependence used in the ideal-model
analysis, but it neither proves FF1 pseudorandomness nor supplies a worst-case
availability guarantee.

\begin{figure}[t]
\centering
\includegraphics[width=\linewidth]{figures/walk_density.pdf}
\caption{Concrete cycle-walk cost across 96 corpus policies supported by the
FF1 backend. Each point is the mean of 32 generation inputs.}
\label{fig:walk-density}
\end{figure}

\subsection{Lifecycle model checking}

The TLA+ specification models \textsc{Active}, \textsc{Selecting},
\textsc{Pending}, \textsc{Submitted}, \textsc{UnknownOutcome}, the four
new/old-password observations, commit, abort, crash, and recovery. With maximum
policy epoch 2 and maximum credential generation 3, TLC 2.19 generated
3,129,888 states, retained 1,006,128 distinct states, and reached depth 56.
TLC reported no violation of: type correctness; no
commit without \textsc{NewOnly}; candidate/committed generation agreement;
stable salt within a lineage; epoch advance on policy change; epoch stability
for same-policy rotation; no committed history reuse; no commit from
\textsc{UnknownOutcome}; and durable pending state before submission.

The recovery-access model checks authorization, public-key and ticket binding,
one-time operation IDs, single-domain factor preservation, ordinary re-sharing,
and Root-Key replacement. TLC generates 6,145,889 states and retains 852,704
distinct states to depth 32. The integrated freshness/compromise model covers
complete epoch/generation rollback, credential-key lineage, Root-Key
replacement, and share-set freshness; it generates 287,151 states and retains
40,292 distinct states to depth 16. TLC reports no violation of the listed
invariants within either configuration.

This is evidence about the configured finite model, not unbounded executions
or every implementation. The model abstracts cryptographic operations as atomic,
does not model target-specific lockout timing, and assumes authenticated local
metadata. The new models additionally abstract signatures, encryption, Shamir,
administrator judgment, and network timing. Their source, configuration, TLC
version, digest, and complete output summaries are included in the artifact.

\section{Security discussion}

The principal unconditional properties are language compliance and
non-repetition conditioned on successful execution of a permutation. Exact
space size is a combinatorial fact. Uniformity and context separation require
ideal-permutation or computational assumptions, respectively. Password
outputs remain limited by $N_P$ and by server-side online/offline attack
conditions; deterministic generation cannot add entropy to a small policy
language.

One exposed password gives one known permutation pair; $q$ exposed passwords
give $q$ pairs. Ideal residual uniformity describes the remaining support, not
forward secrecy and not resistance to concrete key-recovery attacks. A real PRP
replaces ideal equality with computational indistinguishability under query and
domain bounds. The FF1 instance can additionally fail closed at its domain or
walk limits. Uniform history exclusion assumes a fixed authenticated,
deduplicated set rederived before candidate scanning.

Credential-key and Root-Key compromise have different repair scopes. A fresh
credential salt creates an independently derived $K_{cred}$ for one credential;
the first new-lineage password must exclude authenticated history because two
lineages can intersect. A policy-epoch change with the old credential key is
not repair. Root-Key compromise reveals every derivable credential in that
generation. Re-sharing the same value changes shares but not the exposed Root
Key; replacement and remote credential rotation are required.

Permutation injectivity assumes distinct inputs. Rollback can restore a valid
old input and deterministically reproduce its old password without any
collision. The externally anchored multidimensional checkpoint is therefore a
system-level premise for deployment no-repeat. An authenticated local MAC alone
cannot establish global freshness.

Recovery factor preservation is conditional on real administrative
independence. A compromised D obtains $S_D$ and everything callable from its
tokens, but not A signatures by definition. U can be copied and N nodes can
deny service. If one management plane controls D, enough N nodes, and A, the
closure contains two shares and the property no longer holds. Approval reduces
automatic authority; it does not prevent coercion, social engineering, or
malicious authorized administrators.

\section{Limitations and future work}

The policy corpus is historical Web data rather than a current sample of
enterprise deployments. Its 121 exact translations are determined by the
current IR; unsupported PCP semantics are not approximated. One translated
policy reaches the fixed wall-time limit, and the compiler can suffer
product-state explosion. It currently lacks
automaton minimization and optimized multi-pattern matching. Unicode
normalization, dictionary or breached-password services, context-dependent
remote rules, and unbounded semantic predicates are not supported.

The FF1 backend follows a draft minimum-domain rule and bounds cycle walking.
A further 24 compiled corpus policies exceed its 512-bit ceiling. The
enumerable distribution experiment instead uses a test permutation because the
11-word domain is below the FF1 minimum.
A production instantiation should use a reviewed, proof-matched permutation
whose implemented domain and failure semantics exactly match the security
statement, or prove availability bounds acceptable to the application.

The evaluation lacks Windows and Linux measurements, a broader current policy
corpus, comparison on a larger shared policy subset, production target
adapters, fault injection against real services, a recovery usability study,
and an independent cryptographic review of the concrete backend. The TLA+
result is bounded and abstracts cryptography, network timing, lockout, and
storage implementation. The recovery prototype has no distributed transport,
independent enterprise identity integration, human approval study, active
fragment refresh, or Byzantine fragment identification. It measures local
cryptography rather than network or approval latency. D/U/N/A independence is
a deployment requirement rather than a cryptographic fact. These gaps preclude
a production-readiness claim.

A separate theoretical question remains open for this project: dynamic,
intersecting policy languages under adversarial policy changes and a bounded
history state. Before treating that as a new mathematical contribution, it
requires an additional prior-art audit across dynamic FPE, constrained
sampling without replacement, and state lower bounds.

\section{Conclusion}

EPSCD specializes exact finite-language indexing and finite-domain permutation
to password rotation and states the residual credential space after observed
passwords without confusing injectivity with unpredictability. Credential
generation indexes one context-bound permutation; authenticated history,
rekeying, remote evidence, and external freshness preserve its lifecycle
meaning. Recovery reconstructs only the random Root Key and uses closure
analysis to prevent an endpoint share and callable network release from
collapsing the intended threshold. The corpus, published baseline, adversarial
tests, timings, and bounded models establish prototype feasibility but not
production readiness. A proof-matched total finite-domain backend, deployed
recovery transport and authorization, and real target-service evidence remain
the principal next steps.

\appendix
\section{Formal proofs for exposure, history exclusion, and cycle walking}
\label{app:formal-proofs}

\subsection{Residual uniformity after observed credentials}

Let $D$ be a rank domain of size $N$ and let $\pi$ be uniformly sampled from
all permutations on $D$.  Condition on $q<N$ distinct consistent pairs
$(g_i,r_i)$, so $\pi(g_i)=r_i$ for $1\leq i\leq q$.  Let $g^*$ be an
unobserved input and let $r^*$ be any rank outside
$R=\{r_1,\ldots,r_q\}$.

\textbf{Theorem A.1 (observed-credential residual uniformity).}
Under the conditioning above,
\[
\Pr[\pi(g^*)=r^*\mid \pi(g_i)=r_i,1\leq i\leq q]
=\frac{1}{N-q}.
\]
Consequently $\Unrank_P(\pi(g^*))$ is uniform over the $N-q$ legal passwords
whose ranks have not been observed.

\textit{Proof.}
Exactly $(N-q)!$ permutations are consistent with the observed pairs.  Fixing
the additional pair $(g^*,r^*)$ leaves $(N-q-1)!$ consistent permutations.
Their ratio is $1/(N-q)$.  The rank/unrank bijection transfers this statement
to the accepted language. \hfill$\square$

This is standard random-permutation symmetry specialized to the credential
sequence, not a new combinatorial theorem.  For a real finite-domain PRP, let
$\mathcal A$ distinguish the unseen EPSCD rank from this conditional ideal
game after at most $q$ observed or chosen generation queries.  Replacing the
real permutation with a uniformly random one changes $\mathcal A$'s success
probability by at most the PRP advantage at the stated domain, query, and tweak
bounds.  Any preceding replacement of a Root-Key-derived credential key by an
independent key additionally incurs the relevant HKDF/PRF advantage.  The
statement does not turn the ideal equality into an unconditional claim about
the bounded FF1 prototype.

\subsection{Uniform exclusion}

Let $D=\{0,\ldots,N-1\}$ be the rank domain and let
$\Unrank:D\rightarrow\Lang_P$ be a bijection.  The authenticated history is
first rederived, restricted to $\Lang_P$, and deduplicated as a set of strings.
Its inverse rank set is
\[
B=\{\Rank_P(w):w\in E\cap\Lang_P\}\subset D,
\qquad e=|B|.
\]
Thus repeated history descriptors that rederive the same password contribute
one element to $B$, not multiple exclusions.

Fix distinct, unused generation inputs $g_1,g_2,\ldots$ in $D$, and sample
$\pi$ uniformly from the $N!$ permutations of $D$.  Write $Y_i=\pi(g_i)$ and
let
\[
T=\min\{i\geq 1:Y_i\notin B\}.
\]

\textbf{Theorem A.2 (uniform exclusion marginal).}
If $e<N$ and at least $e+1$ distinct generation inputs are available, then
$T\leq e+1$.  For every $z\in D\setminus B$,
\[
\Pr[Y_T=z]=\frac{1}{N-e},
\qquad
\mathbb{E}[T]=\frac{N+1}{N-e+1}.
\]
Consequently, $\Unrank(Y_T)$ is uniform on $\Lang_P\setminus E$.

\textit{Proof.}
For every $t\leq N$ and every tuple of distinct values
$(y_1,\ldots,y_t)$, a uniformly sampled permutation satisfies
\[
\Pr[(Y_1,\ldots,Y_t)=(y_1,\ldots,y_t)]=\frac{1}{(N)_t},
\]
where $(N)_t=N(N-1)\cdots(N-t+1)$.  Hence the images of the distinct scanned
inputs form an ordered sample without replacement.

No more than the $e$ distinct members of $B$ can occur before the first value
outside $B$, which proves $T\leq e+1$.  For a fixed $z\in D\setminus B$,
\[
\Pr[Y_T=z]
=\sum_{t=1}^{e+1}\frac{(e)_{t-1}}{(N)_t}.
\]
The right-hand side does not depend on $z$.  Since termination is certain and
the $N-e$ events $\{Y_T=z\}$ partition the sample space, each probability is
$1/(N-e)$.  Applying the rank/unrank bijection gives the password-space
marginal.

For the expectation, expose a uniformly random ordering of all $N$ ranks and
mark its $N-e$ non-excluded positions.  $T$ is the minimum marked position.
For $0\leq k\leq e$,
\[
\Pr[T>k]=\frac{(e)_k}{(N)_k}
=\frac{\binom{N-k}{N-e}}{\binom{N}{N-e}}.
\]
The tail-sum identity followed by the hockey-stick identity yields
\[
\mathbb{E}[T]
=\sum_{k=0}^{e}\Pr[T>k]
=\frac{1}{\binom{N}{N-e}}
  \sum_{j=N-e}^{N}\binom{j}{N-e}
=\frac{\binom{N+1}{N-e+1}}{\binom{N}{N-e}}
=\frac{N+1}{N-e+1}.
\tag*{$\square$}
\]

The theorem fixes $E$ independently of the newly sampled ideal permutation.
In the protocol, this means the history descriptors are authenticated and
their passwords are rederived before candidate scanning.  Undocumented remote
history and adversarial modification of local descriptors are outside this
probability statement.

\subsection{Cycle-walk termination and ideal tail}

Let $D=\{0,\ldots,N-1\}$ and embed it in
$U=\{0,\ldots,M-1\}$, where $M=2^{\lceil\log_2N\rceil}$.  Starting at
$x\in D$, apply a permutation $P$ on $U$ until its iterate returns to $D$.
Let $W$ count primitive permutation calls, including the accepting call.

\textbf{Lemma A.3 (deterministic termination).}
For every permutation $P$ and every $x\in D$, unbounded cycle walking stops
within $M-N+1$ primitive calls.

\textit{Proof.}
The orbit of $x$ is a finite cycle containing $x\in D$.  Before the next visit
to $D$, it cannot repeat an outside point and there are only $M-N$ such points.
The following call either reaches another member of $D$ or returns to $x$.
\hfill$\square$

\textbf{Lemma A.4 (ideal-permutation tail).}
If $P$ is a uniformly random permutation on $U$, then for every integer
$k\geq0$,
\[
\Pr[W>k]=\frac{(M-N)_k}{(M)_k},
\]
where the numerator is defined as zero for $k>M-N$.  If $N<M$, then
\[
\Pr[W>k]
\leq\left(\frac{M-N}{M}\right)^k<2^{-k}.
\]
If $N=M$, $W=1$ deterministically.

\textit{Proof.}
The event $W>k$ says that the first $k$ orbit images are distinct members of
$U\setminus D$.  Under a random permutation, these images are sampled without
replacement, giving
\[
\frac{M-N}{M}\frac{M-N-1}{M-1}\cdots
\frac{M-N-k+1}{M-k+1}=\frac{(M-N)_k}{(M)_k}.
\]
Each factor is at most $(M-N)/M$.  For a non-power-of-two $N$, the definition
of $M$ gives $M/2<N<M$, hence $(M-N)/M<1/2$.  The power-of-two case has no
outside point and accepts after one call. \hfill$\square$

For a cap of $c$ primitive calls, the ideal-model probability of a cap failure
is $\Pr[W>c]$, not $\Pr[W\geq c]$.  The concrete FF1 implementation is only
computationally related to the ideal permutation experiment.  The bound above
is therefore an ideal-model tail calculation; it is not an unconditional
availability guarantee for FF1 or its implementation.


\section*{Artifact availability}

The artifact contains the Rust implementation, fixed derivation vector,
public-corpus converter and provenance, exact-arithmetic reproduction of the
published Dichopile algorithm, raw JSON measurements, TLA+ model and bounded
checking results, exposure and recovery experiments, and reproduction commands.
Historical internal studies are excluded from the submission artifact. An
anonymized repository URL will be added for review.

\section*{Declaration of competing interest}

The author declares no known competing financial interests or personal
relationships that could have appeared to influence this work.

\section*{Funding}

This research received no specific grant from funding agencies in the public,
commercial, or not-for-profit sectors.

\section*{Data availability}

The evaluation uses the public SOUPS 2022 PCP dataset and records its source
URL and workbook digest. Derived corpus rows, raw measurements, models, and
reproduction scripts are included in the review artifact; an anonymized URL
will be inserted before submission.

\section*{CRediT authorship contribution statement}

Yuanyi Zhang: Conceptualization, Methodology, Software, Validation, Formal
analysis, Investigation, Data curation, Writing--original draft,
Writing--review and editing, Visualization.

\section*{Declaration of generative AI and AI-assisted technologies in the writing process}

During preparation of this work, the author used ChatGPT and Codex to support
language editing, code inspection, and test and documentation drafting. The
author reviewed and edited all resulting material and takes full responsibility
for the content of the publication.

\bibliographystyle{elsarticle-num}
\bibliography{references}

\end{document}
